IMU MinIMU-9 v6 Gyro 
Accéléromètre, gyroscope, magnétomètre

Description et fonctionnement 

Le MinIMU-9 v6 est un module inertiel combinant trois types de capteurs : un accéléromètre, un gyroscope (tous deux intégrés au circuit LSM6DSO), ainsi qu’un magnétomètre (LIS3MDL). Ces trois composants fournissent respectivement l’accélération la vitesse angulaire et le champ magnétique sur 3 axes. Cela permet de calculer l'orientation absolue du capteur dans l’espace et d’estimer les déplacements relatifs. Le MiniIMU-9 v6 que nous utilisons se présente sous la forme d’une carte très compacte. Il communique via une interface I2C, ce qui permet de le connecter facilement à un microcontrôleur, ou bien comme dans notre cas, à la Raspberry Pi 4.

En fonctionnement, l’IMU mesure en continu les mouvements du porteur. L’accéléromètre détecte les accélérations (y compris la gravité), le gyroscope mesure les rotations (vitesse de rotation) et le magnétomètre agit comme une boussole électronique en repérant le nord magnétique. À chaque échantillon, on obtient donc un vecteur d’accélération, un vecteur de vitesse angulaire, et un vecteur de champ magnétique. En utilisant ces données ensemble, on peut estimer l’orientation du capteur à chaque instant (par exemple avec un filtre de fusion de capteurs). Ensuite, en intégrant les accélérations dans un repère corrigé selon cette orientation, on peut essayer de reconstituer la trajectoire suivie.

Contrairement au GPS qui fournit une position absolue mais à faible fréquence, un IMU fournit des mesures relatives à haute fréquence et sans dépendre d’une infrastructure externe, permettant un suivi en temps réel du mouvement, même en intérieur ou dans des tunnels où le GPS. 


Avantage et inconvénients 

Un IMU tel que le MinIMU-9 v6 présente plusieurs atouts dans le contexte du suivi de mouvement. D’abord, il est compact et léger, et s’alimente sur batterie, ce qui le rend facile à porter pendant les acquisitions. Il est également relativement peu coûteux pour la quantité d’informations fournies. De plus, il offre des mesures continues et à haute fréquence de l’activité : il peut suivre en temps réel même des mouvements rapides ou de petites variations de trajectoire. Cela le rend idéal pour
combler les lacunes d’un GPS : l’IMU fonctionne dans tous les environnements, y compris en l’absence totale de signal satellite (espaces couverts). En combinant accéléromètre, gyroscope et magnétomètre, on peut déduire l’orientation du corps et détecter des événements tels que les pas, changements de direction ou inclinaison du terrain. L’IMU apporte donc la possibilité de suivre un parcours sans GPS, ou d’améliorer la précision d’un parcours GPS. Contrairement au GPS qui peut perdre en précision à cause d'obstacles, un IMU reste fiable une fois lancé, car il fonctionne de façon autonome, sans être perturbé par des interférences extérieures. Cela permet de garder un service continu, ce qui est essentiel pour la navigation autonome.

En revanche, les IMU présentent plusieurs limitations. Le principal inconvénient est la dérive au cours du temps : en l’absence de référence externe, de petites erreurs dans les mesures d’accélération ou de rotation s’intègrent et s’accumulent en erreurs de plus en plus grandes sur la vitesse puis la position. En d'autres termes, si l'on utilise seulement un IMU pour estimer une trajectoire, les erreurs vont forcément s'accumuler avec le temps sans ajustements.

En bref, un IMU a plusieurs limites. D'abord, les capteurs qu’il utilise (MEMS) sont sensibles au bruit et à de petites erreurs internes, ce qui fait que la trajectoire estimée dérive peu à peu. De plus, il donne seulement une position relative : sans point de départ connu ou capteur complémentaire (comme un GPS ou une carte), il ne peut pas situer le trajet sur une carte réelle. Il faut aussi régulièrement le calibrer pour corriger les petits défauts sur chaque axe, ce qui peut être long et compliqué. Enfin, le magnétomètre de l’IMU peut être perturbé par des objets métalliques ou des champs magnétiques autour, ce qui fausse l’orientation. En résumé, un IMU est très utile pour suivre un déplacement à court terme ou dans des environnements compliqués, mais il a besoin de corrections pour rester fiable sur la durée.

Sources de bruitage et/ou défauts techniques 
Le MinIMU-9 v6 peut être perturbé par plusieurs types de bruits et défauts liés aux capteurs MEMS qu’il contient. D’abord, chaque capteur (accéléromètre, gyroscope, etc.) génère un bruit électronique naturel, souvent modélisé comme un "bruit blanc". Par exemple, les capteurs du LSM6DSO produisent de petites variations aléatoires dans les mesures d’accélération ou de rotation, ce qui crée de l’incertitude quand on essaie de calculer la position à partir de ces données.					Ensuite, ces capteurs ont aussi des biais (valeurs décalées même quand rien ne bouge) et peuvent dériver dans le temps ou avec la température. Par exemple, un gyroscope peut indiquer qu’il tourne alors qu’il est immobile, et un accéléromètre peut donner une lecture incorrecte de la gravité. Si on ne les calibre pas, ces petites erreurs finissent par s'accumuler et fausser la trajectoire estimée.
Il y a aussi un phénomène appelé "dérive d’intégration" : chaque petite erreur dans les mesures se propage dans les calculs suivants, ce qui entraîne des écarts de plus en plus grands sur la position. C’est pourquoi on doit souvent recourir à des corrections externes (comme des points de référence GPS).
D’autres sources de bruit peuvent affecter les mesures, comme les chocs (ex. : les impacts du pied au sol), les vibrations du boîtier, ou les interférences magnétiques (provoquées par des objets métalliques ou des appareils électroniques proches). Ces perturbations peuvent fausser les données du magnétomètre et rendre difficile l’estimation de l’orientation si elles ne sont pas bien prises en compte.
Tous ces défauts combinés peuvent faire dévier progressivement la trajectoire estimée de la réalité. Pour éviter cela, il est essentiel d’appliquer des méthodes de calibration et de filtrage avancées pour améliorer la fiabilité des données fournies par l’IMU.

Calibration effectuée (démarche, expérimentations et résultats) 

Pour garantir des mesures exploitables, nous avons procédé à une calibration expérimentale de l’IMU MinIMU-9 v6 en le comparant aux capteurs d’un smartphone sur une même séquence de mouvement. L’objectif de cette démarche est de quantifier et corriger les écarts systématiques de l’IMU (biais, facteurs d’échelle) en utilisant le smartphone comme référence. En effet, les smartphones modernes embarquent également des IMUs calibrées en usine et bénéficient souvent d’algorithmes de fusion de capteurs sophistiqués, leurs données de mouvement peuvent donc servir de base de comparaison fiable. Concrètement, nous avons mis le MinIMU-9 et un smartphone l’un à côté de l’autre, puis enregistré simultanément les données des deux dispositifs pendant une série de mouvements (marche en ligne droite, virages, arrêts, accélérations/brusques changements de direction, etc.). Le smartphone enregistrait ses propres données d’accélération et de rotation via une application dédiée(GPS Logger) tandis que la Raspberry Pi enregistrait en parallèle les données brutes de l’IMU externe, le tout avec un référentiel temporel commun.


Choix des paramètres d’acquisition 
Le choix des paramètres de l’IMU, comme la fréquence d’échantillonnage ou la plage de mesure, influence directement la qualité du suivi de trajectoire. Dans notre cas, nous avons opté pour une fréquence de 10 Hz (10 mesures par seconde). Ce rythme permet de suivre globalement la marche d’un utilisateur, avec une cadence typique d’un pas toutes les 0,5 secondes.						  La littérature recommande généralement au moins 20 Hz pour bien capturer les mouvements humains, et nos observations confirment que 10 Hz peut entraîner des erreurs (distance sous-estimée, légers décalages temporels). 			Enfin, même si le GPS fonctionne à 1 Hz, nous avons conservé l’IMU à 10 Hz et utilisé des timestamps pour assurer une bonne synchronisation. Ce réglage offre un bon compromis entre qualité des données, charge de calcul et autonomie du système embarqué.
Filtrage des signaux bruts 
Pour obtenir une trajectoire utilisable à partir des données de l’IMU, il est essentiel de filtrer les signaux bruts afin de réduire le bruit et la dérive. Pour l’instant, nous utilisons un filtrage simple basé sur un modèle physique : à chaque instant, on met à jour la position, la vitesse et l’orientation en supposant un mouvement uniformément accéléré. Cela donne une première estimation du déplacement (phase de prédiction), sans correction par d’autres capteurs.
Cependant, cette méthode reste limitée : les erreurs s’accumulent rapidement en l'absence de recalibrage, notamment à cause du bruit et des biais des capteurs. C’est pourquoi nous prévoyons d’ajouter un filtre de Kalman, qui combinera les données IMU avec d’autres sources comme le GPS. Le Kalman permettra de corriger les dérives et de lisser les données en ajustant en temps réel l’estimation selon la fiabilité de chaque source.
Ce filtre fusionnera par exemple l’estimation rapide du gyroscope avec les repères plus stables fournis par l’accéléromètre, le magnétomètre ou le GPS. À terme, cette approche offrira un suivi de trajectoire plus précis, même en l’absence temporaire de GPS. Nous envisageons aussi d’utiliser un pré-filtrage simple (comme une moyenne glissante) pour éliminer les vibrations, mais un Kalman bien configuré pourra souvent s’en charger directement.




